\documentclass[12pt,a4paper]{article}

% Paquetes
\usepackage[spanish]{babel}
\usepackage[utf8]{inputenc}
\usepackage[T1]{fontenc}
\usepackage{geometry}
\usepackage{setspace}
\usepackage{graphicx}
\usepackage{tikz}
\usetikzlibrary{shapes.geometric, arrows.meta, babel} % Agregamos shapes y babel
\usepackage{enumitem}
\usepackage{titlesec}

\geometry{margin=2.5cm}
\setstretch{1.5}

% Configuración de TikZ
\usetikzlibrary{shapes.geometric, arrows, positioning}

\title{Gestión y Estructura de Procesos Organizacionales\\
\large Enfoque desde la Ingeniería de Sistemas}
\author{}
\date{}

\begin{document}

\maketitle

\tableofcontents
\newpage

\section{Estructura de los Procesos en una Organización}

\subsection{Concepto de Estructura de Procesos}
La estructura de procesos en una organización hace referencia a la forma en que los diferentes procesos se organizan, jerarquizan e interrelacionan para cumplir los objetivos estratégicos. Esta estructura permite comprender el flujo de trabajo desde la toma de decisiones hasta la generación de valor para el cliente.

En ingeniería de procesos, esta organización suele representarse mediante mapas de procesos, los cuales muestran los distintos niveles organizacionales y su interacción.

\subsection{Organización Jerárquica de los Procesos}
Generalmente, los procesos empresariales se estructuran en tres niveles:

\begin{itemize}
    \item \textbf{Procesos estratégicos:} Definen la dirección, planificación y control de la organización.
    \item \textbf{Procesos misionales:} Ejecutan las actividades principales del negocio.
    \item \textbf{Procesos de apoyo:} Proveen recursos y soporte para el funcionamiento organizacional.
\end{itemize}

Estos niveles funcionan como un sistema integrado, no como estructuras aisladas.

\subsection{Relación entre los Tipos de Procesos}
Existe una interdependencia clara entre los distintos tipos de procesos:

\begin{itemize}
    \item Los procesos estratégicos orientan a los misionales.
    \item Los procesos misionales dependen del soporte de los procesos de apoyo.
\end{itemize}

\subsubsection{Ejemplo Aplicado}
Una decisión estratégica como la expansión digital impacta directamente en los procesos de ventas y requiere soporte tecnológico para su implementación.

\subsection{Diagrama de Estructura de Procesos}

\begin{center}
\begin{tikzpicture}[
    node distance=2cm,
    every node/.style={draw, rectangle, rounded corners, align=center, minimum width=3cm, minimum height=1cm}
]

\node (estrategicos) {Procesos Estratégicos};
\node (misionales) [below of=estrategicos] {Procesos Misionales};
\node (apoyo) [below of=misionales] {Procesos de Apoyo};

\draw[->] (estrategicos) -- (misionales);
\draw[->] (apoyo) -- (misionales);

\end{tikzpicture}
\end{center}

\subsection{Conexión entre Procesos}
Los procesos organizacionales están conectados mediante flujos de información, recursos y decisiones. La salida de un proceso puede convertirse en la entrada de otro.

\subsubsection{Ejemplo de Flujo}
\begin{center}
\begin{tikzpicture}[
    node distance=3cm,
    every node/.style={draw, rectangle, rounded corners, align=center}
]

\node (ventas) {Ventas};
\node (logistica) [right of=ventas] {Logística};
\node (finanzas) [right of=logistica] {Finanzas};

\draw[->] (ventas) -- (logistica);
\draw[->] (logistica) -- (finanzas);

\end{tikzpicture}
\end{center}

\subsection{Ejemplo Aplicado: Empresa de Software}
\begin{itemize}
    \item \textbf{Estratégicos:} planificación tecnológica, innovación.
    \item \textbf{Misionales:} desarrollo e implementación de software.
    \item \textbf{Apoyo:} recursos humanos, infraestructura, administración.
\end{itemize}

\section{Segmentación de Procesos}

\subsection{Definición}
La segmentación de procesos consiste en dividir un proceso complejo en partes más pequeñas (subprocesos), facilitando su análisis, control y mejora.

\subsection{Importancia}
La segmentación permite:

\begin{itemize}
    \item Identificar puntos críticos
    \item Detectar cuellos de botella
    \item Mejorar la asignación de responsabilidades
    \item Facilitar el modelamiento en BPMN
\end{itemize}

\subsection{Comprensión del Proceso}
Dividir un proceso permite analizar cada etapa de forma independiente sin perder la visión global, facilitando la documentación y estandarización.

\subsection{Ejemplo Práctico: Gestión de Pedidos}

\begin{center}
\begin{tikzpicture}[
    node distance=2.5cm,
    every node/.style={draw, rectangle, rounded corners, align=center}
]

\node (r1) {Recepción};
\node (r2) [right of=r1] {Validación};
\node (r3) [right of=r2] {Pago};
\node (r4) [right of=r3] {Preparación};
\node (r5) [right of=r4] {Despacho};
\node (r6) [right of=r5] {Entrega};

\draw[->] (r1) -- (r2);
\draw[->] (r2) -- (r3);
\draw[->] (r3) -- (r4);
\draw[->] (r4) -- (r5);
\draw[->] (r5) -- (r6);

\end{tikzpicture}
\end{center}

\section{Procesos, Eficiencia y Mejora Continua}

\subsection{Procesos y Eficiencia Organizacional}
Los procesos bien definidos permiten mejorar la eficiencia organizacional al establecer:

\begin{itemize}
    \item Secuencia de actividades
    \item Uso óptimo de recursos
    \item Reducción de errores
\end{itemize}

\subsection{Estandarización de Procesos}
La estandarización implica definir procedimientos uniformes que aseguren consistencia en los resultados.

\begin{itemize}
    \item Reduce errores humanos
    \item Facilita la capacitación
    \item Mejora la calidad de la información
\end{itemize}

\subsection{Mejora Continua}
La mejora continua busca optimizar procesos de forma constante mediante evaluación y ajustes.

Se basa en el ciclo PDCA:

\begin{center}
\begin{tikzpicture}[
    node distance=2.5cm,
    every node/.style={draw, circle, align=center}
]

\node (plan) {Plan};
\node (do) [right of=plan] {Do};
\node (check) [below of=do] {Check};
\node (act) [left of=check] {Act};

\draw[->] (plan) -- (do);
\draw[->] (do) -- (check);
\draw[->] (check) -- (act);
\draw[->] (act) -- (plan);

\end{tikzpicture}
\end{center}

\subsection{Ejemplo de Mejora Continua}
En un proceso de soporte técnico:

\begin{itemize}
    \item Tiempo inicial: 48 horas
    \item Implementación de sistema de tickets: 36 horas
    \item Estandarización de respuestas: 24 horas
    \item Automatización con IA: 12 horas
\end{itemize}

Esto evidencia la evolución progresiva del proceso mediante mejoras iterativas.

\section{Conclusión}
La gestión de procesos organizacionales permite estructurar el funcionamiento empresarial de manera eficiente, integrada y orientada a resultados. La adecuada organización, segmentación y mejora continua de los procesos constituyen pilares fundamentales en la ingeniería de sistemas.

Estos conceptos son esenciales para el modelamiento en BPMN, ya que facilitan la representación de flujos de trabajo, la identificación de actores y la optimización del desempeño organizacional.

\end{document}
